\honest{Avoiding Your Belief's Real Weak Points}{Eliezer Yudkowsky}{2007}

At my great-grandmother's funeral, my grand-uncle, who had cared for her for years, said through tears that God had called his mother back piece by piece: her memory, then her speech, and last her smile; when God took her smile, he knew she was almost gone.

This puzzled me. In Jewish theology God sustains the universe and chooses every event, but as I understood the rules of religious self-deception, people draw that conclusion only on happy occasions. They say ``God did it'' for a baby girl and not for a stillbirth, and so build a lopsided picture of God's kindness. \nb{Jewish practice points the other way. Jews are required to bless God for bad news as for good, on hearing of a death and, as mourners, at the funeral. Attributing a death to God is prescribed on exactly this occasion.}

Not long afterwards he left the religion. If I had noticed my confusion, I could have predicted this. \nb{The prediction is made after knowing the outcome.} He is the only member of my extended family besides me to have left, as far as I know.

Modern Orthodox Judaism is like no other religion I have heard of: you may notice contradictions, but only in order to explain them away. My example is a rabbi who raises the conflict between the seven days of creation and the 13.7 billion years since the Big Bang, because he has a clever answer involving three Biblical references, a Midrash and a half-understood \textsc{Scientific American} article. The most complicated explanation wins. You attack only targets you know you can defend. \nb{Christians do the same, and have whole organizations devoted to it, such as Reasons to Believe.} You may doubt, but not successfully. \nb{That is true of any group with a shared belief, since the people who doubt successfully leave.}

I then write out the inner thoughts of a typical doubting Orthodox Jew: could the tribes at Sinai have understood the science, did they even have a word for ``billion,'' so the seven days must be a metaphor, with light standing for the Big Bang. I quote no actual Orthodox Jew.

A weaker point is the killing of the firstborn of Egypt, which God does to persuade an unelected Pharaoh to free slaves that God could have simply moved out of the country. Every Orthodox Jew knows the story: the whole Torah is read in synagogue each year, and Passover is named for God passing over Jewish houses during the killing. At the Passover meal, people spill a drop of wine for each of the plagues, to mark the Egyptians' suffering. The rabbis were kinder than the compilers of the Bible and saw the harshness of the plagues, but they stopped short of saying it was wrong. They could afford that much sympathy because science was weaker then. \nb{The custom is medieval, and reading it as sympathy for the Egyptians first appears in writing in 1904, in a lecture by Eduard Baneth.}

Orthodox Jews never even ask whether the killing reflects badly on God, so they need no reply. I supply three replies for them, including ``Murdering babies is okay when God does it!'' Three paragraphs later I grant that there are ``probably a dozen'' replies, and I do not discuss any of them. \nb{A dozen replies means that someone asked.}

My thesis, which I introduce with ``I suspect,'' is that educated believers stay religious because they question their beliefs only where they can defend them, and where rehearsing the standard defense feels strengthening. Answering ``Doesn't Science say that the universe is just meaningless atoms bopping around?'' feels good; it is more comfortable than thinking of an Egyptian mother at her son's crib. Anyone who thinks about an Egyptian mother wailing over her dead son ``is really questioning it, and is probably not going to stay Jewish much longer.'' So a believer who stays was probably never really questioning.

This is not only about Judaism. When people question themselves, they tend to attack strong points that have comforting replies, and they stop at the first reply instead of criticizing it. A better title would be ``Not Spontaneously Thinking About Your Belief's Most Painful Weaknesses.'' I think this is instinct, not training: it hurts, so we avoid it.

My advice: when you doubt a cherished belief, think deliberately about what hurts most. Do not rehearse objections whose standard answers make you feel better. Ask what smart people who disagree would say to your first reply, and then to your second. When you catch yourself flinching from an objection, bring it to the front of your mind. \nb{The post does not apply this advice to its own thesis, that believers stay by avoiding their weak points. The thesis's weakest point is what Jews actually say about the hard passages, and the post does not look at it.}

Then: ``Punch yourself in the solar plexus. Stick a knife in your heart, and wiggle to widen the hole.'' And recite the lines by Eugene Gendlin: what is true is already so; owning up to it does not make it worse; ``People can stand what is true, for they are already enduring it.''
